%% ch04-configuration.tex — JA4proxy Operator's User Guide
%% Chapter 4: Configuration

\chapter{Configuration}
\label{ch:configuration}
\index{configuration}

\newthought{All runtime behaviour is controlled by \code{config/proxy.yml}.}
This chapter walks through every significant configuration section, explains what
each setting does, and gives guidance on choosing values for your environment.

The configuration file uses YAML. Every key is documented with an inline comment
in the file itself; this chapter provides the operational context that comments
cannot.

\section{Hot Reload}
\label{sec:hotreload}
\index{hot reload}
\index{SIGHUP}

Configuration changes take effect without restarting the proxy. Send \code{SIGHUP}
to the proxy process:

\begin{bashcode}
kill -HUP $(docker compose exec ja4proxy cat /var/run/proxy.pid)
\end{bashcode}

The proxy logs a confirmation:
\begin{lstlisting}[language=bash, numbers=none]
INFO | config | event=reload | status=ok | keys_changed=3
\end{lstlisting}

The Management UI (when present) triggers reload via a Redis pub/sub message —
the proxy subscribes to a reload channel and applies new configuration on receipt.

Configuration sections that \emph{cannot} be hot-reloaded and require a full
restart: listen port, Redis URL, TLS certificate paths, and the number of worker
processes. These are logged as \code{restart\_required} in the reload response.

\section{The Dial}
\index{dial}
\index{dial!rollout}

\begin{yamlcode}
# Blocking aggression: 0 = monitor only, 100 = full blocking
# Default: 0 (never block on first deploy)
dial: 0
\end{yamlcode}

The dial is the single most consequential configuration setting. Get it wrong and
you either block legitimate users or provide no protection.

\subsection{Dial Rollout Procedure}

Follow this progression. Spend at least 24--48 hours at each stage before proceeding,
and watch the Grafana action distribution and false-positive indicators throughout.

\begin{table}[ht]
  \small
  \begin{tabularx}{\linewidth}{llX}
    \toprule
    \textbf{Dial} & \textbf{Effect} & \textbf{What to watch for} \\
    \midrule
    0 & Monitor only & Baseline. All traffic passes. Review score distribution. \\
    25 & Flag only & High-scoring traffic is flagged in metrics. No blocking. \\
    50 & Rate limiting active & Automated clients start hitting rate limits. Check
         for false positives in rate-limited IPs. \\
    75 & Tarpit + rate limit & High-risk connections enter tarpit. Verify no
         legitimate API clients are affected. \\
    100 & Full blocking & Block threshold applies. Last check: review BLOCKED log
         lines for any unexpected legitimate traffic. \\
    \bottomrule
  \end{tabularx}
  \caption{Dial rollout stages}
\end{table}

\begin{dangerbox}[Do not jump to dial=100]
The dial increment limit exists for a reason. A jump from 0 to 100 on a production
system, without first reviewing scores at intermediate stages, risks immediately
blocking legitimate users. Always progress gradually.
\end{dangerbox}

\subsection{Dial and Bypass Interaction}

ALLOW bypasses (browser ALPN, JA4 whitelist, mTLS) are \emph{unaffected by the dial}.
They always allow. Disabling a bypass is the correct way to route that category
through the scorer; the dial then controls whether scoring results in a block.

BLOCK bypasses (JA4 blacklist, country block, Spamhaus) also operate independently
of the dial at their default setting — they produce hard blocks before scoring.
When a BLOCK bypass is disabled, those connections enter the scorer, and the dial
then controls what happens to them.

\section{Country Filtering}
\index{country filtering}
\index{GeoIP}

\begin{yamlcode}
# Country whitelist: only these countries are allowed
# Requires GeoIP database (make update-geoip)
country_whitelist_enabled: false
country_whitelist:
  - IE
  - GB
  - US

# Country blacklist: these countries are always blocked
# Takes effect only when country_whitelist_enabled is false
country_blacklist_enabled: false
country_blacklist:
  - KP
  - RU
  - CN
\end{yamlcode}

Whitelist and blacklist are mutually exclusive: if \configkey{country\_whitelist\_enabled}
is \code{true}, the blacklist is ignored (the whitelist already excludes everything not
listed). If both are disabled, country filtering has no effect.

\begin{warningbox}[Whitelist vs blacklist]
The whitelist is the more robust option for deployments where you know your legitimate
audience is geographically bounded. The blacklist is more convenient but less effective
— attackers are not constrained to operate from the countries you have blocked.

VPN usage means country codes are unreliable. Do not rely on country filtering as
your primary defence.
\end{warningbox}

\section{JA4 Fingerprint Lists}
\index{JA4!blacklist}
\index{JA4!whitelist}
\index{fingerprint lists}

\begin{yamlcode}
security:
  # Exact-match whitelist: these fingerprints are always allowed (bypass scoring)
  whitelist:
    - t13d1113h2_8daaf6152771_b0da82dd1658  # Chrome 120 on Windows
    - t13d1113h2_8daaf6152771_02713d6af862  # Chrome 120 on macOS

  # Pattern-match whitelist: regex patterns matched against fingerprint strings
  # Use for version-resilient matching when the full fingerprint changes with browser updates
  whitelist_patterns:
    - "t13d111[0-9]h2_.*"  # Any Chrome-like TLS 1.3 fingerprint

  # Exact-match blacklist: these fingerprints are immediately blocked (hard block)
  blacklist:
    - t13d0912_a1b2c3d4e5f6_987654321abc  # Sliver C2
    - t12d1302_c0ffee123456_aabbccddeeff  # Cobalt Strike default

  # Human-readable labels for known fingerprints (used in logs and dashboard)
  fingerprint_labels:
    t13d1113h2_8daaf6152771_b0da82dd1658: "Chrome 120 (Windows)"
    t13d1113h2_8daaf6152771_02713d6af862: "Chrome 120 (macOS)"
    t13d0912_a1b2c3d4e5f6_987654321abc: "Sliver C2"
\end{yamlcode}

\subsection{Managing the Whitelist}

Add fingerprints to the whitelist when:
\begin{itemize}
  \item A legitimate API client or monitoring tool is being scored too harshly
  \item An internal service has a non-browser TLS fingerprint you know is safe
  \item A specific customer's tooling is repeatedly triggering rate limits
\end{itemize}

Use \configkey{whitelist\_patterns} for browser fingerprints that change with each
browser update, to avoid constant maintenance. The pattern \code{t13d111[0-9]h2\_.*}
matches all TLS 1.3 connections with Chrome-like cipher ordering, regardless of
the specific extension hashes.

\begin{notebox}[Whitelist patterns and performance]
Patterns are evaluated as regular expressions for each connection. Keep the pattern
list short (under 20 entries) and anchor patterns where possible to avoid
backtracking. Exact matches in \configkey{whitelist} are O(1) hash lookups.
\end{notebox}

\subsection{Managing the Blacklist}

Add fingerprints to the blacklist when you have high confidence that a fingerprint
represents a malicious tool and that no legitimate software shares it. The JA4
community database at \url{https://ja4db.com} is a useful reference for known
tool fingerprints.

The blacklist produces an immediate RST (hard block) before scoring. This is the
most efficient way to drop known-bad clients.

\begin{warningbox}[Shared fingerprints]
Some penetration testing tools use the same underlying TLS library as legitimate
software. Before blacklisting a fingerprint, verify it is not shared with any
legitimate client in your environment. Fingerprints from generic HTTP libraries
(like Python's \code{requests} or Node.js's \code{https}) are often shared.
\end{warningbox}

\section{Rate Limiting}
\index{rate limiting}
\index{rate limiting!strategies}

\begin{yamlcode}
rate_limiting:
  enabled: true

  # Three independent strategies — any, all, or a majority can trigger
  rate_limit_strategies:
    by_ip:
      enabled: true
      limit: 60          # connections per window
      window_seconds: 60
      action: ban        # ban, block, or tarpit

    by_ja4:
      enabled: true
      limit: 200         # same fingerprint across all IPs
      window_seconds: 60
      action: tarpit

    by_ip_ja4_pair:
      enabled: true
      limit: 30          # same IP + same fingerprint
      window_seconds: 60
      action: ban

  # How to combine strategy decisions when multiple trigger
  # majority: trigger if >=2 of 3 strategies say yes (default)
  # any: trigger if any strategy says yes (aggressive)
  # all: trigger if all enabled strategies say yes (conservative)
  multi_strategy_policy: majority
\end{yamlcode}

\subsection{Strategy Selection}

The three strategies catch different attack patterns:

\begin{description}[leftmargin=4.5cm]
  \item[\configkey{by\_ip}] Catches any single source sending high volume,
    regardless of tool. The bluntest instrument — catches bots, scanners,
    and DDoS sources. Low false-positive risk for normal users.

  \item[\configkey{by\_ja4}] Catches distributed attacks using the same tool
    across many source IPs. A botnet of 1000 IPs each sending 2 connections
    might not trigger the IP strategy, but will trigger the JA4 strategy if
    the fingerprint is consistent across the botnet.

  \item[\configkey{by\_ip\_ja4\_pair}] The most precise strategy. Catches
    a single source using a single tool intensively, while allowing the same
    IP to use multiple different tools (e.g., a security researcher running
    multiple scanners from one IP).
\end{description}

\subsection{Multi-Strategy Policy}

\configkey{multi\_strategy\_policy: majority} is the default and recommended
setting. It requires two of the three strategies to agree before triggering
an action. This reduces false positives from any single strategy misfiring,
while still catching distributed attacks that trigger multiple strategies.

Use \configkey{any} only if you need maximum sensitivity and are comfortable
with a higher false-positive rate. Use \configkey{all} only if you have
persistent false positive issues with the majority policy.

\section{Bypass Conditions}
\index{bypass conditions}
\index{security policy}

\begin{yamlcode}
security_policy:
  # ALLOW bypasses — short-circuit scoring
  alpn_browser_bypass:
    enabled: true        # h2/h1 ALPN → always allowed

  ja4_whitelist_bypass:
    enabled: true        # JA4 in whitelist → always allowed

  mtls_bypass:
    enabled: true        # Valid mTLS client cert → always allowed

  static_ip_allowlist:
    enabled: true        # IP in static allowlist → always allowed

  # BLOCK bypasses — hard block before scoring
  ja4_blacklist_bypass:
    enabled: true        # JA4 in blacklist → immediate RST

  country_blacklist_bypass:
    enabled: true        # Country in blocklist → immediate block

  spamhaus_bypass:
    enabled: true        # Spamhaus DROP/EDROP match → immediate block

  tls_version_bypass:
    enabled: true        # TLS 1.0/1.1 → immediate RST
\end{yamlcode}

Every bypass condition is independently configurable. The defaults are sensible
for most deployments; here is when you would change them:

\begin{table}[ht]
  \small
  \begin{tabularx}{\linewidth}{lX}
    \toprule
    \textbf{Bypass} & \textbf{When to disable} \\
    \midrule
    \configkey{alpn\_browser\_bypass} & Almost never. Disabling this means legitimate
      browsers can be blocked, which is the primary false-positive risk. Only disable
      if scoring specific h2 API clients is critical to your use case. \\
    \midrule
    \configkey{ja4\_whitelist\_bypass} & If you want whitelisted fingerprints to be
      scored for observability purposes but still allowed. \\
    \midrule
    \configkey{mtls\_bypass} & If you want to score mTLS clients rather than
      unconditionally trust them. \\
    \midrule
    \configkey{spamhaus\_bypass} & If you want to investigate traffic from Spamhaus-listed
      ranges rather than hard-blocking it. Disabling routes those connections through
      the scorer with a +80 risk signal contribution. \\
    \midrule
    \configkey{tls\_version\_bypass} & If you need to score rather than immediately
      reject TLS 1.0/1.1 clients — e.g., to quantify how much legacy TLS your
      infrastructure actually receives. \\
    \bottomrule
  \end{tabularx}
  \caption{When to disable bypass conditions}
\end{table}

\begin{warningbox}[Startup warnings for disabled bypasses]
If you disable any high-risk bypass, the proxy emits a \code{WARN} log at startup
and increments a Prometheus counter. This is intentional — it is a reminder that
a non-default bypass configuration is active. The warning format:
\begin{lstlisting}[language=bash, numbers=none]
WARN | policy | event=bypass_disabled | bypass=alpn_browser_bypass |
       effect=browser traffic will be scored; false positive risk elevated
\end{lstlisting}
\end{warningbox}

\section{Additional Configuration Sections}

\subsection{Signal Weights}
\index{risk scorer!weights}

Individual signal weights can be tuned in the \configkey{scorer} section. The
defaults are calibrated from empirical testing; only adjust them if you have
clear evidence that a specific signal is over- or under-contributing in your
environment. \emph{See the Technical Reference Manual for signal weight documentation.}

\subsection{Beaconing Detection}
\index{beaconing}

\begin{yamlcode}
beaconing:
  enabled: true
  min_connections: 5       # minimum connections before analysis
  cv_threshold: 0.3        # coefficient of variation threshold
  window_seconds: 3600     # look-back window
\end{yamlcode}

The beaconing detector looks for connections from the same IP+JA4 pair at
suspiciously regular intervals. A coefficient of variation below the threshold
(very regular timing) triggers a risk signal. The default threshold of 0.3
catches C2 beaconing without flagging APIs with regular polling intervals.

\subsection{AbuseIPDB}
\index{AbuseIPDB}

\begin{yamlcode}
abuseipdb:
  enabled: true            # set to false if no API key
  api_key: ""              # or set via ABUSEIPDB_API_KEY env var
  score_threshold: 50      # minimum score to generate a risk signal
  cache_ttl_seconds: 3600  # cache successful lookups for 1 hour
  hard_block_threshold: 0  # 0 = never hard-block via AbuseIPDB alone
\end{yamlcode}

AbuseIPDB score is used as one signal among many. The \configkey{hard\_block\_threshold}
defaults to 0 (disabled) because AbuseIPDB scores reflect community reports that may
include shared IPs (NAT, VPNs). A high AbuseIPDB score contributes to the composite
risk score, but does not by itself produce a hard block.

\subsection{Redis Connection}
\index{Redis!configuration}

\begin{yamlcode}
redis:
  host: redis              # hostname or socket path
  port: 6379
  db: 0
  # unix_socket: /var/run/redis/redis.sock  # 30-40% lower latency than TCP
\end{yamlcode}

The Unix domain socket option is significantly faster than TCP for same-host Redis.
If your proxy and Redis run on the same host, configure the socket path.
